\chapter*{阅读指引}

邮件分类器在测试集上表现良好，是否就可以自动处理用户的全部邮件？实际邮件可能换了
模板，不同错误可能造成不同损失，某些用户可能更容易受影响，助手也可能泄露内容或执行越权
操作。测试成绩只回答了其中一部分问题。本部分从这些具体要求出发，检查模型与整个系统能否
胜任实际任务，以及失败后怎样限制损害。

《可信机器学习基础》先建立共同的判断方法：明确谁在什么情境中使用系统、希望它怎样
行动，再检查理论、测试和运行记录能支持哪些结论，并安排监测、停用与恢复。后续各章沿着
使用中的疑问展开：预测概率是否值得依赖，环境变化后是否仍能正确判断，不同群体是否得到
公平待遇，行为是否符合实际目标与权限，解释能否帮助理解和补救，私人信息是否受到保护，
主动攻击者又能否突破这些防线。

实际系统往往同时涉及多项要求。一个模型的平均错误率下降，并不说明它已经更公平、更
安全或更尊重隐私。阅读各章时，都要检查所研究的具体问题、采用的证据和结论成立的条件，
再把这些结果放回同一个系统，判断它能承担哪些任务、在哪些情形下仍需人工复核或限制功能。
